\chapter{Estado del Arte del Tema}
\pendiente{fundamentación teórica (LSA, reconocimiento de señas con visión por computadora, LSTM sobre landmarks) con citas académicas.}

\section{Contexto actual en Latinoamérica, Argentina y Tucumán}
\pendiente{población sorda, Ley 27.710 de reconocimiento de la LSA, situación en las guardias.}

\section{Innovación tecnológica del proyecto}
\pendiente{qué aporta el proyecto frente al estado del arte.}

\section{Impacto social del proyecto}
\pendiente{beneficiarios, de qué forma se benefician y otros impactos.}

\section{Necesidades y oportunidades que justifican el proyecto}
\pendiente{situación actual contra situación esperada.}

\section{Beneficios cuantitativos y cualitativos}
\pendiente{beneficios medibles y no medibles.}

\section{Tecnologías similares}
\pendiente{comparativa de traductores de lengua de señas existentes y justificación de la elección.}
